% =============================================================================
%  SECTIONS RECOMMENDED FOR REWRITING — BY THE AUTHOR IN CROATIAN
%  These texts show the strongest "translation-speak" patterns.
%  Rewrite them in Croatian (your original language), then re-translate
%  to English with an AI tool or editor.
% =============================================================================

%%% ======================================================================= %%%
%%%   SECTION 4 (04 peft.tex) — Survey taxonomy and landscape paragraph      %%%
%%% ======================================================================= %%%
%
% ORIGINAL (lines 225-228):
%
% The PEFT taxonomy is not merely an academic classification exercise.
% In the context of a P2P adapter marketplace, it becomes a concrete
% design foundation with real engineering consequences. LoRA is the
% natural choice for transmission efficiency — its parameter delta is
% the smallest of any reviewed method, and its compact matrix
% representation is the most practical to move across an unpredictable
% P2P network with variable bandwidth and latency. Bottleneck adapters,
% on the other hand, present the stronger case for composability — their
% modularity is precisely the property that composition architectures
% such as AdapterFusion exploit, and it is not easily replicated by
% reparameterisation methods.
%
% These two advantages point in opposite directions, and it is not
% possible to fully optimise for both simultaneously. The choice between
% them is therefore a tradeoff that the proposed system must navigate
% explicitly — not a conclusion that the literature resolves cleanly.
% The proof-of-concept design follows the direction of LoRA as the
% default adapter format, primarily because of its prevalence in the
% reviewed literature, its demonstrated adaptability across model scales,
% and its influence on the federated and serving systems that form the
% closest prior work. But the composability argument for bottleneck
% adapters does not disappear — it is set aside for pragmatic reasons,
% and remains an open design question for future iterations of the
% system.
%
% WHY REWRITE: This paragraph alternates perfectly balanced "on one hand
% ... on the other hand" constructions that are characteristic of AI
% summarisation — too balanced, no commitment. A human author writing
% from the Croatian original would naturally take a position more
% directly.
%
% === WHAT I WOULD SAY IN CROATIAN (example) ===
% U ovom radu koristim LoRA kao primarni adapter format za POC, ali ne
% zato što je LoRA objektivno bolja od bottleneck adapters za ovaj
% slučaj. Praktični razlozi su presudili: većina pregledane literature
% koristi LoRA, skalabilna je, a najveći utjecaj na kasnije sustave je
% upravo kroz LoRA. Bottleneck adapters nisu lošiji — oni su jednako dobri,
% samo nisu praktični za ovaj setup. Ovo je svjestan trade-off, ne
% akademska neodlučnost.
%%% ======================================================================= %%%

\subsection{Survey Taxonomy and Positioning}
\label{sec:peft:positioning_reframe}

[Rewrite from Croatian here → replace the original paragraph 04_peft.tex ll.225-228]


%%% ======================================================================= %%%
%%%   SECTION 5 (05 adapter composition.tex) — positioning paragraph       %%%
%%% ======================================================================= %%%
%
% ORIGINAL (lines 179-186):
%
% The composition methods reviewed in this chapter presuppose that a
% collection of adapters already exists in a single, centralised  182
% location. Composition and adapter discovery are often treated as
% separate concerns, but in a P2P network they are tightly coupled — a
% peer cannot compose what it has not yet found. The obstacle is not
% whether composition works technically; LoraHub and LoraRetriever
% demonstrate that it does. The real obstacle is that the substrate on
% which these methods depend does not exist in P2P environments.
% No mechanism currently exists for constructing a compact
% description of an adapter's capabilities without access to its
% original training data — a gap that directly motivates contributions
% $\star$R1 and $\star$R2. No mechanism exists for routing a query
% toward capable adapters in the absence of a central index or
% directory — the gap that motivates $\star$A1. And even if all of
% these pieces were in place, the existing composition methods remain
% blind to the realities of distributed network operation: peer churn,
% variable latency, and bandwidth constraints. Addressing this
% infrastructure layer is the motivation behind $\star$S1 and $\star$S2.
%
% WHY REWRITE: Overly patterned "No mechanism exists... no mechanism
% exists... and even if all of these pieces were in place..." is
% formulaic. A Croatian original would likely lead with the main
% claim and use simpler coordination, less parallelism.
%
% === WHAT I WOULD SAY IN CROATIAN (example) ===
% Metode prikazane u ovoj sekciji sve pretpostavljaju centralizirani
% setup gdje su adapteri već prikupljeni na jednom mjestu. U
% peer-to-peer kontekstu problem nije to što te metode ne rade — rade,
% LoraHub i LoraRetriever to dokazuju. Problem je u tome što u P2P
% okruženju ne postoji ta centralna kolekcija adaptera koju bi te
% metode mogle koristiti. Zato su potrebni novi mehanizmi: adapters
% capabilities representation ($\star$R1), behavioural fingerprints
% ($\star$R2), DHT discovery ($\star$A1), i P2P infrastruktura
% ($\star$S1, $\star$S2).
%%% ======================================================================= %%%

\subsection{Implications for Decentralised Deployment}
\label{sec:composition:implications_reframe}

[Rewrite from Croatian here → replace the original paragraph 05_adapter_composition.tex ll.179-186]


%%% ======================================================================= %%%
%%%   SECTION 6 (06 inference systems.tex) — multi-tenant gap paragraph     %%%
%%% ======================================================================= %%%
%
% ORIGINAL (lines 175-176):
%
% The existing serving literature addresses the multi-tenant inference
% problem with considerable technical depth — but it leaves one
% fundamental question unanswered: where does the adapter catalogue
% come from, and who controls it? Every system reviewed in this
% chapter inherits its adapter collection from a centralised upstream
% process that is never designed, only assumed.
%
% WHY REWRITE: The "where does the catalogue come from" question is a
% good insight, but the framing ("never designed, only assumed") reads
% as AI-generated rhetorical closure. The Croatian original need only
% state: "ovi sustavi pretpostavljaju da adapteri negdje postoje i da
% je njihov katalog centraliziran. Problem je u tome što u P2P
% izostaje svaki centralni izvor adaptera — i nijedan sustav o tome
% ni ne raspravlja."
%%% ======================================================================= %%%

\subsection{Gap: Missing Source of Adapters}
\label{sec:inference:gap_reframe}

[Rewrite from Croatian here → replace original paragraph 06_inference_systems.tex ll.175-176]


%%% ======================================================================= %%%
%%%   SECTION 7 (07 p2p federated.tex) — subsection ll.221-226             %%%
%%% ======================================================================= %%%
%
% ORIGINAL (lines 221-226):
%
% Federated PEFT contributes three properties that are directly
% relevant to the target architecture: it reduces communication
% overhead to the adapter level, which in practice yields bandwidth
% savings of up to 90\%; it provides formal guarantees of data
% privacy and isolation; and it tolerates heterogeneous client
% resources. These are substantial contributions. Yet federated PEFT
% carries one limitation that is insurmountable in a P2P adapter
% serving context: the central server remains the sole coordination
% point for aggregation, distribution, and synchronisation of adapters
% across participants.
%
% P2P distributed learning addresses precisely that limitation. By
% combining DHT-based lookup for locating information, gossip-based
% propagation for disseminating it, and convergence guarantees under
% peer churn, P2P systems provide the infrastructure for removing
% the central server entirely. What P2P systems lack, however, is
% adapter semantics — there are no task-specific modules, no
% capability matching, and no composition logic anywhere in this
% literature.
%
% This is not a failure of either field. Each body of work is
% capable of complementing the other's blind spots — federated PEFT
% solves what P2P ignores, and P2P solves what federated PEFT cannot
% relinquish. What is missing is a concrete protocol that bridges the
% two: one that can discover a useful adapter across the network,
% retrieve it from whichever peer currently holds it, perform local
% composition, and serve the inference result — all without any
% single point of control or a single source of truth.
%
% WHY REWRITE: The "This is not a failure... Each body... capable of
% complementing..." conclusion is structurally very AI-like in
% its hedging balance. This is your key argument — it would be
% stronger stated directly: "Ni jedan ni drugi pristup nije u krivu
% — oba su potrebna, ali razdvojena. Federated PEFT rješava ono što
% P2P ne pokriva, i obrnuto. Ono što nedostaje je konkretni protokol
% koji spaja ta dva svijeta."
%%% ======================================================================= %%%

\subsection{Bridging Gap: Federated vs P2P}
\label{sec:p2p:gap_reframe}

[Rewrite from Croatian here → replace the corresponding part of 08_p2p_federated.tex (P2P/federated section; formerly 07_p2p_federated.tex)]

